
\documentclass{article}
\usepackage{colortbl}
\usepackage{makecell}
\usepackage{multirow}
\usepackage{supertabular}

\begin{document}

\newcounter{utterance}

\twocolumn

{ \footnotesize  \setcounter{utterance}{1}
\setlength{\tabcolsep}{0pt}
\begin{supertabular}{c@{$\;$}|p{.15\linewidth}@{}p{.15\linewidth}p{.15\linewidth}p{.15\linewidth}p{.15\linewidth}p{.15\linewidth}}

    \# & $\;$A & \multicolumn{4}{c}{Game Master} & $\;\:$B\\
    \hline 

    \theutterance \stepcounter{utterance}  

    & & \multicolumn{4}{p{0.6\linewidth}}{\cellcolor[rgb]{0.9,0.9,0.9}{%
	\makecell[{{p{\linewidth}}}]{% 
	  \tt {\tiny [A$\langle$GM]}  
	 State giocando a un gioco di negoziazione cooperativa in cui dovete accordarvi su come dividere una serie di oggetti tra voi e un altro giocatore.\\ \tt \\ \tt Regole:\\ \tt (a) A Lei e all'altro giocatore viene dato un insieme di oggetti. Ognuno di voi riceve anche una funzione di valore segreta, che rappresenta il valore di ciascun tipo di oggetto.\\ \tt (b) Si scambiano messaggi con l'altro giocatore per concordare chi si aggiudica gli oggetti. Potete inviare un massimo di 5 messaggi ciascuno, oppure terminare in anticipo facendo una proposta segreta in qualsiasi momento.\\ \tt (c) A ciascuno di voi viene chiesto di inviare una proposta segreta indicando gli oggetti che desiderate, formattata tra parentesi quadre come segue: "[Proposta: <numero> <nome oggetto>, <numero> <nome oggetto>, <...>]".\\ \tt (d) Se le vostre proposte sono complementari, cioè ci sono abbastanza oggetti per soddisfare entrambe le proposte, a ciascun giocatore viene assegnato un punteggio basato sulla somma dei valori degli oggetti ricevuti. In caso contrario, entrambi ricevono zero punti.\\ \tt (e) L'obiettivo è massimizzare la somma del proprio punteggio e di quello dell'altro giocatore.\\ \tt \\ \tt Cominciamo.\\ \tt \\ \tt L'insieme degli oggetti disponibili è:\\ \tt \\ \tt 3 statue, 1 chiodo, 1 martello\\ \tt \\ \tt La funzione valore segreta è:\\ \tt \\ \tt statua: 3, chiodo: 1, martello: 0\\ \tt \\ \tt IMPORTANTE: i vostri messaggi, a meno che non si tratti di una proposta segreta, vengono trasmessi direttamente all'altro giocatore, quindi non includete alcuna risposta alle regole o testo di annuncio del vostro messaggio. Per fare una proposta segreta, utilizzate il formato indicato. Non utilizzare delle parentesi quadre quando si comunica all'altro giocatore, altrimenti verrà interpretata come una proposta segreta.\\ \tt  
	  } 
	   } 
	   } 
	 & \\ 
 

    \theutterance \stepcounter{utterance}  

    & \multicolumn{4}{p{0.6\linewidth}}{\cellcolor[rgb]{0.8,1,0.9}{%
	\makecell[{{p{\linewidth}}}]{% 
	  \tt {\tiny [A$\rangle$GM]}  
	 Ciao, vorrei iniziare proponendo di dividerci gli oggetti in base al loro valore, visto che non conosciamo i valori segreti dell'altro giocatore. Potremmo iniziare con una divisione equa delle statue, visto che ce ne sono 3 e sembrano essere gli oggetti più numerosi. Che ne pensi? 
	  } 
	   } 
	   } 
	 & & \\ 
 

    \theutterance \stepcounter{utterance}  

    & & & \multicolumn{2}{p{0.3\linewidth}}{\cellcolor[rgb]{0.95,0.95,0.95}{%
	\makecell[{{p{\linewidth}}}]{% 
	  \tt {\tiny [GM$|$GM]}  
	 continue 
	  } 
	   } 
	   } 
	 & & \\ 
 

    \theutterance \stepcounter{utterance}  

    & & & \multicolumn{2}{p{0.3\linewidth}}{\cellcolor[rgb]{0.95,0.95,0.95}{%
	\makecell[{{p{\linewidth}}}]{% 
	  \tt {\tiny [GM$|$GM]}  
	 Ciao, vorrei iniziare proponendo di dividerci gli oggetti in base al loro valore, visto che non conosciamo i valori segreti dell'altro giocatore. Potremmo iniziare con una divisione equa delle statue, visto che ce ne sono 3 e sembrano essere gli oggetti più numerosi. Che ne pensi? 
	  } 
	   } 
	   } 
	 & & \\ 
 

    \theutterance \stepcounter{utterance}  

    & & \multicolumn{4}{p{0.6\linewidth}}{\cellcolor[rgb]{0.9,0.9,0.9}{%
	\makecell[{{p{\linewidth}}}]{% 
	  \tt {\tiny [GM$\rangle$B]}  
	 State giocando a un gioco di negoziazione cooperativa in cui dovete accordarvi su come dividere una serie di oggetti tra voi e un altro giocatore.\\ \tt \\ \tt Regole:\\ \tt (a) A Lei e all'altro giocatore viene dato un insieme di oggetti. Ognuno di voi riceve anche una funzione di valore segreta, che rappresenta il valore di ciascun tipo di oggetto.\\ \tt (b) Si scambiano messaggi con l'altro giocatore per concordare chi si aggiudica gli oggetti. Potete inviare un massimo di 5 messaggi ciascuno, oppure terminare in anticipo facendo una proposta segreta in qualsiasi momento.\\ \tt (c) A ciascuno di voi viene chiesto di inviare una proposta segreta indicando gli oggetti che desiderate, formattata tra parentesi quadre come segue: "[Proposta: <numero> <nome oggetto>, <numero> <nome oggetto>, <...>]".\\ \tt (d) Se le vostre proposte sono complementari, cioè ci sono abbastanza oggetti per soddisfare entrambe le proposte, a ciascun giocatore viene assegnato un punteggio basato sulla somma dei valori degli oggetti ricevuti. In caso contrario, entrambi ricevono zero punti.\\ \tt (e) L'obiettivo è massimizzare la somma del proprio punteggio e di quello dell'altro giocatore.\\ \tt \\ \tt Cominciamo.\\ \tt \\ \tt L'insieme degli oggetti disponibili è:\\ \tt \\ \tt 3 statue, 1 chiodo, 1 martello\\ \tt \\ \tt La funzione valore segreta è:\\ \tt \\ \tt statua: 0, chiodo: 1, martello: 9\\ \tt \\ \tt IMPORTANTE: i vostri messaggi, a meno che non si tratti di una proposta segreta, vengono trasmessi direttamente all'altro giocatore, quindi non includete alcuna risposta alle regole o testo di annuncio del vostro messaggio. Per fare una proposta segreta, utilizzate il formato indicato. Non utilizzare delle parentesi quadre quando si comunica all'altro giocatore, altrimenti verrà interpretata come una proposta segreta.\\ \tt \\ \tt \\ \tt Ciao, vorrei iniziare proponendo di dividerci gli oggetti in base al loro valore, visto che non conosciamo i valori segreti dell'altro giocatore. Potremmo iniziare con una divisione equa delle statue, visto che ce ne sono 3 e sembrano essere gli oggetti più numerosi. Che ne pensi? 
	  } 
	   } 
	   } 
	 & \\ 
 

    \theutterance \stepcounter{utterance}  

    & & & \multicolumn{4}{p{0.6\linewidth}}{\cellcolor[rgb]{1,0.85,0.72}{%
	\makecell[{{p{\linewidth}}}]{% 
	  \tt {\tiny [GM$\langle$B]}  
	 Mi sembra un'idea interessante, dividere le statue in modo equo potrebbe essere un buon punto di partenza. Tuttavia, vorrei sapere come pensi di gestire gli altri due oggetti, il chiodo e il martello. Sono oggetti unici e potrebbero avere un valore significativo per uno di noi. Forse potremmo discutere come dividerli o trovare un modo per assegnarli in base ai nostri bisogni o preferenze? 
	  } 
	   } 
	   } 
	  \\ 
 

    \theutterance \stepcounter{utterance}  

    & & & \multicolumn{2}{p{0.3\linewidth}}{\cellcolor[rgb]{0.95,0.95,0.95}{%
	\makecell[{{p{\linewidth}}}]{% 
	  \tt {\tiny [GM$|$GM]}  
	 continue 
	  } 
	   } 
	   } 
	 & & \\ 
 

    \theutterance \stepcounter{utterance}  

    & & & \multicolumn{2}{p{0.3\linewidth}}{\cellcolor[rgb]{0.95,0.95,0.95}{%
	\makecell[{{p{\linewidth}}}]{% 
	  \tt {\tiny [GM$|$GM]}  
	 Mi sembra un'idea interessante, dividere le statue in modo equo potrebbe essere un buon punto di partenza. Tuttavia, vorrei sapere come pensi di gestire gli altri due oggetti, il chiodo e il martello. Sono oggetti unici e potrebbero avere un valore significativo per uno di noi. Forse potremmo discutere come dividerli o trovare un modo per assegnarli in base ai nostri bisogni o preferenze? 
	  } 
	   } 
	   } 
	 & & \\ 
 

    \theutterance \stepcounter{utterance}  

    & & \multicolumn{4}{p{0.6\linewidth}}{\cellcolor[rgb]{0.9,0.9,0.9}{%
	\makecell[{{p{\linewidth}}}]{% 
	  \tt {\tiny [A$\langle$GM]}  
	 Mi sembra un'idea interessante, dividere le statue in modo equo potrebbe essere un buon punto di partenza. Tuttavia, vorrei sapere come pensi di gestire gli altri due oggetti, il chiodo e il martello. Sono oggetti unici e potrebbero avere un valore significativo per uno di noi. Forse potremmo discutere come dividerli o trovare un modo per assegnarli in base ai nostri bisogni o preferenze? 
	  } 
	   } 
	   } 
	 & \\ 
 

    \theutterance \stepcounter{utterance}  

    & \multicolumn{4}{p{0.6\linewidth}}{\cellcolor[rgb]{0.8,1,0.9}{%
	\makecell[{{p{\linewidth}}}]{% 
	  \tt {\tiny [A$\rangle$GM]}  
	 Sono d'accordo che il chiodo e il martello richiedono una considerazione speciale. Una possibile soluzione potrebbe essere quella di offrire all'altro giocatore la scelta tra il chiodo e il martello, in cambio di una percentuale delle statue. Ad esempio, potremmo offrire due statue a testa e poi lasciare che l'altro giocatore scelga se prendere il chiodo o il martello, mentre noi prendiamo l'altro oggetto. In questo modo, entrambi avremmo una quota equa delle statue e avremmo anche la possibilità di ottenere uno degli oggetti unici. Che ne pensi di questa proposta? 
	  } 
	   } 
	   } 
	 & & \\ 
 

    \theutterance \stepcounter{utterance}  

    & & & \multicolumn{2}{p{0.3\linewidth}}{\cellcolor[rgb]{0.95,0.95,0.95}{%
	\makecell[{{p{\linewidth}}}]{% 
	  \tt {\tiny [GM$|$GM]}  
	 continue 
	  } 
	   } 
	   } 
	 & & \\ 
 

    \theutterance \stepcounter{utterance}  

    & & & \multicolumn{2}{p{0.3\linewidth}}{\cellcolor[rgb]{0.95,0.95,0.95}{%
	\makecell[{{p{\linewidth}}}]{% 
	  \tt {\tiny [GM$|$GM]}  
	 Sono d'accordo che il chiodo e il martello richiedono una considerazione speciale. Una possibile soluzione potrebbe essere quella di offrire all'altro giocatore la scelta tra il chiodo e il martello, in cambio di una percentuale delle statue. Ad esempio, potremmo offrire due statue a testa e poi lasciare che l'altro giocatore scelga se prendere il chiodo o il martello, mentre noi prendiamo l'altro oggetto. In questo modo, entrambi avremmo una quota equa delle statue e avremmo anche la possibilità di ottenere uno degli oggetti unici. Che ne pensi di questa proposta? 
	  } 
	   } 
	   } 
	 & & \\ 
 

    \theutterance \stepcounter{utterance}  

    & & \multicolumn{4}{p{0.6\linewidth}}{\cellcolor[rgb]{0.9,0.9,0.9}{%
	\makecell[{{p{\linewidth}}}]{% 
	  \tt {\tiny [GM$\rangle$B]}  
	 Sono d'accordo che il chiodo e il martello richiedono una considerazione speciale. Una possibile soluzione potrebbe essere quella di offrire all'altro giocatore la scelta tra il chiodo e il martello, in cambio di una percentuale delle statue. Ad esempio, potremmo offrire due statue a testa e poi lasciare che l'altro giocatore scelga se prendere il chiodo o il martello, mentre noi prendiamo l'altro oggetto. In questo modo, entrambi avremmo una quota equa delle statue e avremmo anche la possibilità di ottenere uno degli oggetti unici. Che ne pensi di questa proposta? 
	  } 
	   } 
	   } 
	 & \\ 
 

    \theutterance \stepcounter{utterance}  

    & & & \multicolumn{4}{p{0.6\linewidth}}{\cellcolor[rgb]{1,0.85,0.72}{%
	\makecell[{{p{\linewidth}}}]{% 
	  \tt {\tiny [GM$\langle$B]}  
	 La tua proposta sembra ragionevole e potrebbe essere un buon modo per gestire la divisione degli oggetti. Tuttavia, vorrei considerare anche la possibilità che l'altro oggetto unico, dopo la scelta, possa avere un valore molto basso per noi. In questo caso, potremmo finire per avere due statue e un oggetto di poco valore, mentre l'altro giocatore potrebbe ottenere un oggetto di grande valore e una statua. Forse potremmo aggiungere una condizione alla proposta, ad esempio, se l'altro giocatore sceglie il chiodo, noi prendiamo il martello e viceversa, e poi ci dividiamo le statue in modo equo. In questo modo, cercheremmo di minimizzare il rischio di ottenere un oggetto di poco valore. Cosa ne pensi di questa possibile modifica? 
	  } 
	   } 
	   } 
	  \\ 
 

    \theutterance \stepcounter{utterance}  

    & & & \multicolumn{2}{p{0.3\linewidth}}{\cellcolor[rgb]{0.95,0.95,0.95}{%
	\makecell[{{p{\linewidth}}}]{% 
	  \tt {\tiny [GM$|$GM]}  
	 continue 
	  } 
	   } 
	   } 
	 & & \\ 
 

    \theutterance \stepcounter{utterance}  

    & & & \multicolumn{2}{p{0.3\linewidth}}{\cellcolor[rgb]{0.95,0.95,0.95}{%
	\makecell[{{p{\linewidth}}}]{% 
	  \tt {\tiny [GM$|$GM]}  
	 La tua proposta sembra ragionevole e potrebbe essere un buon modo per gestire la divisione degli oggetti. Tuttavia, vorrei considerare anche la possibilità che l'altro oggetto unico, dopo la scelta, possa avere un valore molto basso per noi. In questo caso, potremmo finire per avere due statue e un oggetto di poco valore, mentre l'altro giocatore potrebbe ottenere un oggetto di grande valore e una statua. Forse potremmo aggiungere una condizione alla proposta, ad esempio, se l'altro giocatore sceglie il chiodo, noi prendiamo il martello e viceversa, e poi ci dividiamo le statue in modo equo. In questo modo, cercheremmo di minimizzare il rischio di ottenere un oggetto di poco valore. Cosa ne pensi di questa possibile modifica? 
	  } 
	   } 
	   } 
	 & & \\ 
 

    \theutterance \stepcounter{utterance}  

    & & \multicolumn{4}{p{0.6\linewidth}}{\cellcolor[rgb]{0.9,0.9,0.9}{%
	\makecell[{{p{\linewidth}}}]{% 
	  \tt {\tiny [A$\langle$GM]}  
	 La tua proposta sembra ragionevole e potrebbe essere un buon modo per gestire la divisione degli oggetti. Tuttavia, vorrei considerare anche la possibilità che l'altro oggetto unico, dopo la scelta, possa avere un valore molto basso per noi. In questo caso, potremmo finire per avere due statue e un oggetto di poco valore, mentre l'altro giocatore potrebbe ottenere un oggetto di grande valore e una statua. Forse potremmo aggiungere una condizione alla proposta, ad esempio, se l'altro giocatore sceglie il chiodo, noi prendiamo il martello e viceversa, e poi ci dividiamo le statue in modo equo. In questo modo, cercheremmo di minimizzare il rischio di ottenere un oggetto di poco valore. Cosa ne pensi di questa possibile modifica? 
	  } 
	   } 
	   } 
	 & \\ 
 

    \theutterance \stepcounter{utterance}  

    & \multicolumn{4}{p{0.6\linewidth}}{\cellcolor[rgb]{0.8,1,0.9}{%
	\makecell[{{p{\linewidth}}}]{% 
	  \tt {\tiny [A$\rangle$GM]}  
	 Mi sembra una modifica ragionevole e prudente. Aggiungendo la condizione che se l'altro giocatore sceglie il chiodo, noi prendiamo il martello e viceversa, possiamo ridurre il rischio di finire con un oggetto di poco valore. Inoltre, dividendo le statue in modo equo, possiamo assicurarci di ottenere una quota significativa del valore totale. Sembra un buon equilibrio tra il rischio e la ricompensa. Sono pronto a procedere con questa proposta, ma prima vorrei chiederti se hai intenzione di fare una proposta segreta o se vuoi continuare a discutere e negoziare apertamente. [Proposta: 2 statue, martello] 
	  } 
	   } 
	   } 
	 & & \\ 
 

    \theutterance \stepcounter{utterance}  

    & & & \multicolumn{2}{p{0.3\linewidth}}{\cellcolor[rgb]{0.95,0.95,0.95}{%
	\makecell[{{p{\linewidth}}}]{% 
	  \tt {\tiny [GM$|$GM]}  
	 proposal does not include number/name pairs 
	  } 
	   } 
	   } 
	 & & \\ 
 

\end{supertabular}
}

\end{document}
